% Módulo de expansión para capitulos/marco/04-piezo.tex.
%
% Integración prevista: después de la forma débil y antes de la subsección
% «Antecedente local y lectura de sus figuras». El módulo aplica las
% convenciones ya definidas en el capítulo al caso P5 y no repite la
% metodología de entrenamiento. Usa las figuras que el equipo preparará en
% figuras/revision2026/; ambas son esquemas, no soluciones FEM.

\subsection{Caso concreto: bimorfo efectivo de PVDF}
\label{subsec:piezo-bimorfo-expansion}

El caso P5 aplica el constitutivo piezoeléctrico ya definido a un bimorfo
efectivo de PVDF. Dos capas unidas reciben una diferencia de potencial entre
sus caras principales. La cadena física que interesa distinguir puede resumirse
en dos líneas:
\[
  \begin{aligned}
  \text{potencial }\varphi
    &\;\longrightarrow\;\text{campo }\mathbf E=-\nabla\varphi
    \;\longrightarrow\;\text{acoplamiento }e,\\[-.2ex]
  &\text{deformación y tensión}
    \;\longrightarrow\;\text{flexión}.
  \end{aligned}
\]
El efecto directo relaciona deformación con respuesta eléctrica; el inverso
relaciona campo eléctrico con deformación. En la viga, cada capa tendería a
cambiar su longitud según su polarización, pero la unión impide que las dos
capas alcancen deformaciones longitudinales independientes. Esa incompatibilidad
es la causa mecánica de la curvatura \cite{khan2018bimorph,lahmer2008}.

\subsubsection{Geometría, variables y orientación}

El dominio efectivo es
\[
  \Omega=(0,L)\times(0,H),\qquad
  L=0.1\,\mathrm{m},\qquad H=0.001\,\mathrm{m},
\]
con $x$ a lo largo de la viga y $y$ en su altura. La interfaz perfectamente
unida es $y=H/2$; cada capa tiene espesor $H/2$. Se escriben
\[
  \Omega^-=(0,L)\times(0,H/2),\qquad
  \Omega^+=(0,L)\times(H/2,H),
\]
y se toma como normal común $\boldsymbol n=(0,1)$, dirigida de la capa
inferior a la superior. Las fronteras son
\[
  \Gamma_L=\{0\}\times[0,H],\quad
  \Gamma_R=\{L\}\times[0,H],\quad
  \Gamma_B=[0,L]\times\{0\},\quad
  \Gamma_T=[0,L]\times\{H\}.
\]

Las incógnitas son el desplazamiento material
$\mathbf d=(d_x,d_y)$, en metros, y el potencial $\varphi$, en voltios. El
desplazamiento eléctrico $\mathbf D$, en $\mathrm{C}/\mathrm{m}^{2}$, es un
flujo electromagnético y no un desplazamiento material. Se conserva la
convención del capítulo: la deformación usa cortante ingenieril y
\[
  \mathbf E=-\nabla\varphi=(-\varphi_{,x},-\varphi_{,y}).
\]
La ilustración siguiente identifica materiales y condiciones; la gráfica de solución se obtiene por separado mediante el FEM acoplado.


\begin{figure}[htbp]\centering
\includegraphics[width=\linewidth,height=.30\textheight,keepaspectratio]{p5-problem.pdf}
\caption[P5: ilustración del bimorfo piezoeléctrico]{Ilustración de P5 con Fiziko: dos capas de PVDF adheridas, empotramiento izquierdo, electrodos a 100 y 0 V, e interfaz material. El campo apunta hacia menor potencial. La geometría se exagera para distinguir capas y condiciones; la respuesta mecánica se obtiene mediante el FEM acoplado.}\label{fig:p5-bimorfo}
\end{figure}


\subsubsection{Constitutivo particular y ecuaciones fuertes}

La forma constitutiva general del capítulo se mantiene. Para hacer explícita
la convención de P5, sea
\[
  e_{21}^{0}=0.02904\,\mathrm{C}/\mathrm{m}^{2},\qquad
  e_{22}^{0}=0.05157\,\mathrm{C}/\mathrm{m}^{2},\qquad
  \eta_-=+1,\quad \eta_+=-1,
\]
y defínase, en cada capa,
\begin{equation}
  e_\ell=\eta_\ell
  \begin{bmatrix}
    0&0&0\\
    e_{21}^{0}&e_{22}^{0}&0
  \end{bmatrix},
  \qquad \ell\in\{-,+\}.
  \label{eq:p5-expansion-e-sign}
\end{equation}
Esta notación reproduce los signos opuestos de $e_{21}$ y $e_{22}$ de la
tabla del caso, y fija en cero los otros elementos de $e$. Si
$\mathbf d=(u,v)$ y los coeficientes son constantes dentro de cada material,
la ley del capítulo se lee
\begin{align}
  \sigma_{xx}&=C_{11}u_{,x}+C_{12}v_{,y}
       +\eta_\ell e_{21}^{0}\varphi_{,y},
  &\sigma_{yy}&=C_{12}u_{,x}+C_{22}v_{,y}
       +\eta_\ell e_{22}^{0}\varphi_{,y},
  \label{eq:p5-expansion-stress-normal}\\
  \sigma_{xy}&=C_{66}(u_{,y}+v_{,x}),
  &D_x&=-\kappa_x\varphi_{,x},
  \label{eq:p5-expansion-stress-shear}\\
  &&D_y&=\eta_\ell(e_{21}^{0}u_{,x}+e_{22}^{0}v_{,y})
       -\kappa_y\varphi_{,y}.
  \label{eq:p5-expansion-flux}
\end{align}
Los términos con $e$ son los que transmiten el campo eléctrico al esfuerzo y
la deformación al flujo. El término $-\kappa_y\varphi_{,y}$ es la respuesta
dieléctrica; el signo negativo procede de $\mathbf E=-\nabla\varphi$ y no de
una permitividad negativa.

Sin fuerza volumétrica ni carga libre, el equilibrio cuasiestático en cada
capa es
\begin{equation}
  \partial_x\sigma_{xx}+\partial_y\sigma_{xy}=0,\qquad
  \partial_x\sigma_{xy}+\partial_y\sigma_{yy}=0,\qquad
  \partial_xD_x+\partial_yD_y=0.
  \label{eq:p5-expansion-balances}
\end{equation}
Al sustituir las expresiones anteriores, las tres PDE fuertes por material
quedan
\begin{align}
  C_{11}u_{,xx}+C_{66}u_{,yy}+(C_{12}+C_{66})v_{,xy}
      +\eta_\ell e_{21}^{0}\varphi_{,xy}&=0,\\
  C_{66}v_{,xx}+C_{22}v_{,yy}+(C_{12}+C_{66})u_{,xy}
      +\eta_\ell e_{22}^{0}\varphi_{,yy}&=0,\\
  \eta_\ell e_{21}^{0}u_{,xy}+\eta_\ell e_{22}^{0}v_{,yy}
      -\kappa_x\varphi_{,xx}-\kappa_y\varphi_{,yy}&=0.
  \label{eq:p5-expansion-pdes}
\end{align}
Estas formas se aplican dentro de $\Omega^-$ y $\Omega^+$ por separado. No
se atraviesa $y=H/2$ con una sola PDE de coeficientes suaves: el cambio de
material se completa con transmisión de campos y flujos.

\subsubsection{Fronteras e interfaz}

El contrato mecánico y eléctrico del caso concreto es
\begin{align}
  \mathbf d&=\mathbf 0 &&\text{en }\Gamma_L, &
  \sigma\boldsymbol n&=\mathbf 0
       &&\text{en }\Gamma_R\cup\Gamma_T\cup\Gamma_B,\\
  \varphi&=100\,\mathrm{V} &&\text{en }\Gamma_T, &
  \varphi&=0 &&\text{en }\Gamma_B,\\
  \mathbf D\cdot\boldsymbol n&=0
       &&\text{en }\Gamma_L\cup\Gamma_R.
  \label{eq:p5-expansion-bc}
\end{align}
En los extremos aislados, $D_x=0$. En las caras con potencial fijado no se
prescribe simultáneamente el flujo: la carga superficial del electrodo es una
salida del problema. En las esquinas se usan las trazas de cada segmento, sin
imponer una normal única que combine datos incompatibles en un punto aislado.

Sea $[q]=q^+-q^-$ el salto al pasar de la capa inferior a la superior. Una
interfaz perfectamente unida, sin carga superficial, exige
\begin{equation}
  [\mathbf d]=\mathbf 0,\qquad [\varphi]=0,\qquad
  [\sigma\boldsymbol n]=\mathbf 0,\qquad
  [\mathbf D\cdot\boldsymbol n]=0
  \quad\text{en }y=H/2.
  \label{eq:p5-expansion-interface}
\end{equation}
Con $\boldsymbol n=(0,1)$, las condiciones naturales equivalen a
$[\sigma_{xy}]=[\sigma_{yy}]=0$ y $[D_y]=0$. Esto no exige continuidad de
todo el tensor ni de todas las derivadas. En general no se impone una
condición independiente sobre el flujo tangencial; en este P5 concreto, sin
embargo, la primera fila de $e$ es cero y $\kappa_x$ es igual en las capas.
Si la traza de $\varphi$ es diferenciable a lo largo de la interfaz, entonces
$\varphi_{,x}$ y $D_x=-\kappa_x\varphi_{,x}$ también son continuos. Las
derivadas normales de $u$, $v$ y $\varphi$ sí pueden saltar aunque los tres
campos sean continuos.

\begin{figure}[htbp]
  \centering
  \includegraphics[width=.88\linewidth]{p5-interfaz.pdf}
  \caption[Continuidad de potencial y flujo en una interfaz dieléctrica.]{Ejemplo
  didáctico escalar, normalizado y sin acoplamiento piezoeléctrico. En
  $\hat y=y/H$, con $\hat\kappa_1=1$, $\hat\kappa_2=2$,
  $\hat\varphi(0)=0$ y $\hat\varphi(1)=1$, las pendientes son $4/3$ y
  $2/3$, mientras $\hat D=-4/3$ es continuo en $\hat y=1/2$. La figura
  ilustra continuidad de campo y flujo frente a salto de derivada; no es una
  solución del PVDF P5 ni usa sus permitividades.}
  \label{fig:p5-interfaz}
\end{figure}

Para leer la figura sin atribuirle física adicional, considérese la ley
escalar
\begin{equation}
  \hat D=-\hat\kappa(\hat y)\frac{\mathrm d\hat\varphi}{\mathrm d\hat y},
  \qquad
  \hat\kappa_1=1\ (0<\hat y<1/2),\qquad
  \hat\kappa_2=2\ (1/2<\hat y<1).
  \label{eq:p5-expansion-didactic-interface}
\end{equation}
Sin carga libre, $\hat D$ es constante. Las pendientes
$\mathrm d\hat\varphi/\mathrm d\hat y=4/3$ y $2/3$ producen
$\hat D=-4/3$ a ambos lados, y
$(1/2)(4/3)+(1/2)(2/3)=1$ satisface el voltaje total normalizado. El ejemplo
solo explica una regla de transmisión; no asigna $\hat\kappa_1$ o
$\hat\kappa_2$ al PVDF y no es un perfil extraído de FEM.

\subsubsection{Lectura de signos y distinción de esquema y solución}

Si el potencial aumenta desde la cara inferior a la superior,
$\varphi_{,y}>0$ y $E_y=-\varphi_{,y}<0$. Para una capa libre de tensión,
la tendencia constitutiva es
\begin{equation}
  s_\ell^{\,\mathrm{libre}}=-(C^E)^{-1}e_\ell^T\nabla\varphi.
  \label{eq:p5-expansion-free-strain}
\end{equation}
Con los signos opuestos de $e_{21}$ y $e_{22}$ y los valores positivos de
$C_{11}$, $C_{22}$ y $C_{12}$ de la especificación, la lectura axial da
$\varepsilon_{xx}^{\mathrm{libre},+}>0$ arriba y
$\varepsilon_{xx}^{\mathrm{libre},-}<0$ abajo en esta convención. La unión
convierte esa diferencia en una tendencia de curvatura hacia abajo, que es lo
que predice esta lectura constitutiva. La flecha y la distribución de curvatura se obtienen del campo acoplado; la solución numérica se muestra al final de esta subsección.
Como control de escala, si se sustituye el gradiente nominal
$\varphi_{,y}=E_0=10^5\,\mathrm{V}/\mathrm{m}$ en una capa libre y se usan
los $C$ y $e$ especificados, la capa superior da aproximadamente
$(\varepsilon_{xx},\varepsilon_{yy})=(7.04\times10^{-7},
2.16\times10^{-6})$ y la inferior cambia ambos signos. Es una deducción
constitutiva local, no una solución FEM de la viga.

El signo de $D_y$ no se deduce solo de $\varphi_{,y}$ porque en
\eqref{eq:p5-expansion-flux} compiten la parte dieléctrica y la polarización
inducida por la deformación. Invertir el voltaje debe invertir las respuestas
del modelo lineal; reducir 100 V a 50 V debe reducir aproximadamente a la
mitad sus magnitudes después de verificar la solución. Son comprobaciones de
simetría, no mediciones ya realizadas.

La referencia numérica P5 se resolvió en el rectángulo efectivo
de dimensiones \(L=0.1\,\mathrm{m}\) y \(H=10^{-3}\,\mathrm{m}\), dividido en
dos capas con interfaz alineada en \(y=H/2\). Se usaron triángulos
conformes cuadráticos \(P_2\), una regla de cuadratura de Dunavant de siete
puntos y el sistema acoplado físico
\[
\begin{bmatrix}K_{dd}&K_{d\phi}\\K_{\phi d}&-K_{\phi\phi}\end{bmatrix}
\begin{bmatrix}d\\\phi\end{bmatrix}=0,
\]
con el orden global \((d_x,d_y,\phi)\) por nodo. El borde izquierdo se fijó
en \(d=(0,0)\); los demás bordes mecánicos son naturales, mientras que
\(\phi=0\) en la cara inferior, \(\phi=100\,\mathrm{V}\) en la superior y los
extremos laterales tienen \(D\!\cdot n=0\). La interfaz comparte los grados
de libertad conformes y por ello impone continuidad de \(d\) y \(\phi\), con
las condiciones de transmisión naturales para tracción y flujo normal.

La malla fina usada en la figura tiene 3072 elementos y 6369 nodos \(P_2\).
La flecha en \((L,H/2)\) es
\(d_y=-1.059806\times10^{-5}\,\mathrm{m}\); el cambio relativo entre las dos mallas más finas es
0.311\%. La carga libre del electrodo superior, por unidad de espesor
fuera del plano, es
\(Q_T=1.053905\times10^{-6}\,\mathrm{C/m}\) y cambia
0.002\% en esa misma comparación. El residual infinito de los grados
de libertad libres es
\(2.284210\times10^{-18}\) en las unidades escaladas de P5. Para no invertir la convención
de signos, \(Q_T=-\int_{\Gamma_T}D\!\cdot n\,ds\): el residual eléctrico
de potencial prescrito es el flujo saliente del material y tiene el signo
opuesto al de la carga libre sobre el electrodo. La deformada del panel
correspondiente se amplifica por el factor calculado
28.3071; ese factor es solo visual. Esta salida es una referencia
numérica didáctica de la formulación efectiva 2-D, no una validación
experimental ni una aprobación de resultados de investigación.

\begin{figure}[htbp]\centering
\includegraphics[width=\linewidth]{p5-solution.pdf}
\caption[P5: solución FEM piezoeléctrica acoplada]{Solución FEM acoplada calculada para P5 con elementos triangulares
cuadráticos \(P_2\) y la malla fina alineada con la interfaz. (a) Campo
vertical \(d_y\) en micrómetros; (b) potencial \(\phi\) en voltios; (c)
contorno exterior e interfaz originales y deformados, cuya escala visual es el factor calculado
28.3071; (d) cambio de la flecha de punta y de la carga libre del
electrodo superior al refinar. La deformada se dibuja con coordenadas
geométricas \(x,y\) en metros y con razón de aspecto ajustada para facilitar
la lectura. Como \(L/H=100\), los ejes no usan la misma escala gráfica; además,
el desplazamiento se multiplica por el factor indicado solo para ser visible.
La carga mostrada usa
\(Q_T=-\int_{\Gamma_T}D\!\cdot n\,ds\), mientras el flujo saliente del
material tiene el signo contrario.
}\label{fig:p5-solution}
\end{figure}
